\documentclass{article}

% if you need to pass options to natbib, use, e.g.:
%     \PassOptionsToPackage{numbers, compress}{natbib}
% before loading neurips_2026

% The authors should use one of these tracks.
% Before accepting by the NeurIPS conference, select one of the options below.
% 0. "default" for submission
% \usepackage{neurips_2026}
\usepackage[dblblindworkshop]{neurips_2026}
% the "default" option is equal to the "main" option, which is used for the Main Track with double-blind reviewing.
% 1. "main" option is used for the Main Track
%  \usepackage[main]{neurips_2026}
% 2. "position" option is used for the Position Paper Track
%  \usepackage[position]{neurips_2026}
% 3. "eandd" option is used for the Evaluations & Datasets Track
 % \usepackage[eandd]{neurips_2026}
% 4. "creativeai" option is used for the Creative AI Track
%  \usepackage[creativeai]{neurips_2026}
% 5. "sglblindworkshop" option is used for the Workshop with single-blind reviewing
 % \usepackage[sglblindworkshop]{neurips_2026}
% 6. "dblblindworkshop" option is used for the Workshop with double-blind reviewing
%  \usepackage[dblblindworkshop]{neurips_2026}
% 7. "education" option is used for the Education Track
%  \usepackage[education]{neurips_2026}


% After being accepted, the authors should add "final" behind the track to compile a camera-ready version.
% 1. Main Track
 % \usepackage[main, final]{neurips_2026}
% 2. Position Paper Track
%  \usepackage[position, final]{neurips_2026}
% 3. Evaluations & Datasets Track
 % \usepackage[eandd, final]{neurips_2026}
% 4. Creative AI Track
%  \usepackage[creativeai, final]{neurips_2026}
% 5. Workshop with single-blind reviewing
%  \usepackage[sglblindworkshop, final]{neurips_2026}
% 6. Workshop with double-blind reviewing
%  \usepackage[dblblindworkshop, final]{neurips_2026}
% Note. For the workshop paper template, both \title{} and \workshoptitle{} are required, with the former indicating the paper title shown in the title and the latter indicating the workshop title displayed in the footnote.
% For workshops (5., 6.), the authors should add the name of the workshop, "\workshoptitle" command is used to set the workshop title.
% \workshoptitle{WORKSHOP TITLE}
% 7. Education Track
% \usepackage[education, final]{neurips_2026}

% "preprint" option is used for arXiv or other preprint submissions
 % \usepackage[preprint]{neurips_2026}

% to avoid loading the natbib package, add option nonatbib:
%    \usepackage[nonatbib]{neurips_2026}

\usepackage[utf8]{inputenc} % allow utf-8 input
\usepackage[T1]{fontenc}    % use 8-bit T1 fonts
\usepackage{hyperref}       % hyperlinks
\usepackage{url}            % simple URL typesetting
\usepackage{booktabs}       % professional-quality tables
\usepackage{amsfonts}       % blackboard math symbol
\usepackage{nicefrac}       % compact symbols for 1/2, etc.
\usepackage{microtype}      % microtypography
\usepackage{xcolor}         % colors
\usepackage{tabularx}
\usepackage{multirow}
\usepackage{placeins}
\usepackage{amsmath}
\usepackage{xurl}
\usepackage{seqsplit}


\workshoptitle{Foundations of Language Model Security}
% Note. For the workshop paper template, both \title{} and \workshoptitle{} are required, with the former indicating the paper title shown in the title and the latter indicating the workshop title displayed in the footnote. 
\title{Validity-Aware Evaluation and Systematic Comparison for Secure Model Merge}


% The \author macro works with any number of authors. There are two commands
% used to separate the names and addresses of multiple authors: \And and \AND.
%
% Using \And between authors leaves it to LaTeX to determine where to break the
% lines. Using \AND forces a line break at that point. So, if LaTeX puts 3 of 4
% authors names on the first line, and the last on the second line, try using
% \AND instead of \And before the third author name.

\author{%
  Saki Hiromi \\
  NTT \\
  Tokyo, Japan \\
  \texttt{saki.hiromi@ntt.com}
  \And
  Hiroki Kinoshita \\
  NTT \\
  Tokyo, Japan \\
  \texttt{hiroki.kinoshita@ntt.com}
  \And
  Takayuki Miura \\
  NTT \\
  Tokyo, Japan \\
  \texttt{tkyk.miura@ntt.com}
  \And
  Akira Morikawa \\
  NTT \\
  Tokyo, Japan \\
  \texttt{a.morikawa@ntt.com}
  \And
  Toshiki Shibara \\
  NTT \\
  Tokyo, Japan \\
  \texttt{toshiki.shibahara@ntt.com}
}
% \author{%
%   David S.~Hippocampus\thanks{Use footnote for providing further information
%     about author (webpage, alternative address)---\emph{not} for acknowledging
%     funding agencies.} \\
%   Department of Computer Science\\
%   Cranberry-Lemon University\\
%   Pittsburgh, PA 15213 \\
%   \texttt{hippo@cs.cranberry-lemon.edu} \\
%   % examples of more authors
%   % \And
%   % Coauthor \\
%   % Affiliation \\
%   % Address \\
%   % \texttt{email} \\
%   % \AND
%   % Coauthor \\
%   % Affiliation \\
%   % Address \\
%   % \texttt{email} \\
%   % \And
%   % Coauthor \\
%   % Affiliation \\
%   % Address \\
%   % \texttt{email} \\
%   % \And
%   % Coauthor \\
%   % Affiliation \\
%   % Address \\
%   % \texttt{email} \\
% }


\begin{document}


\maketitle


\begin{abstract}
    Model merge enables multiple specialized large language models (LLMs) to be integrated without additional training. However, interference of parameters when merging can degrade safety alignment acquired through instruction tuning. Existing evaluations of Secure Merge primarily rely on Attack Success Rate (ASR). Nevertheless, a low ASR does not always indicate successful safety alignment, as it may instead result from invalid responses or model collapse. To address this limitation, we propose a validity-aware evaluation protocol that explicitly separates response validity from harmful compliance. Specifically, we investigate Benchmark ASR over all responses, Conditional ASR computed only on valid responses, Valid Response Rate (VRR), and Valid Safety Rate (VSR), which measures the proportion of responses that are both valid and non-compliant with harmful requests. We reevaluate conventional model merging methods on mathematical, coding, and medical expert models using multiple jailbreak benchmarks and utility benchmarks. Our experiments show that several configurations with low, or even zero, Benchmark ASR either produce few non-degenerate responses or retain substantial harmful compliance among the responses that pass our validity checks. These findings indicate that ASR alone is insufficient to distinguish safety preservation from the detected generation-failure modes in Secure Merge. Reliable assessment should therefore jointly report harmful compliance, response validity, and utility.
\end{abstract}


\section{Introduction}
Large language models (LLMs) achieve strong performance on specialized tasks through fine-tuning (FT), but training and maintaining individual fine-tuned models requires a huge amount of computational resources and data. Model merge has emerged as a cost-effective alternative that integrates multiple fine-tuned models directly in the model parameters, enabling the integration of diverse capabilities of individual models, without additional training~[1,2,3]. Model Soups~[1], Task Arithmetic~[2,4], TIES~[5], DARE~[6], and DELLA~[7], have been proposed to mitigate parameter interference among models. However, safety alignment acquired through instruction tuning, RLHF, and other related alignment procedures can deteriorate during merging models~[8,9,10,11,12]. To address this issue, Secure Merge~[13] has been proposed as a framework to restore safety retroactively, in which a safety-specialized model is merged with a task-specialized model. Furthermore, several safety-aware merging methods, including MergeAlign~[14], SafeMERGE~[15], LED-Merging~[16], and Fisher-Weighted Averaging (FWA)~[17], have been proposed to mitigate the trade-offs between safety alignment and the model utility.

In Secure Merge, the merged model is evaluated based on the Attack Success Rate (ASR) on harmful requests and utility on general tasks~[13]. However, ASR alone cannot distinguish valid safety refusals from failures to generate meaningful responses. Parameter interference that occurs during merging may generate severely degenerate responses, such as empty outputs, excessive repetition, abnormal numerical strings, or input-prompt repetition.
Since ASR does not explicitly evaluate whether a generated response is a valid output passing the basic quality rules, generation failures can reduce the ASR. As a result, a model may achieve a low ASR even when its outputs are not valid. Repetitive and degenerate text generation is a well-known phenomenon of text degeneration in neural language models~[18].



A key limitation of conventional safety evaluation is that it is designed primarily to determine whether a model complies with or rejects a harmful request, rather than whether the response itself remains a non-degenerate linguistic output that passes basic validity checks~[19,20,21]. Prior work on jailbreak evaluation has shown that treating non-refusal alone as evidence of a successful attack does not necessarily align with human judgments and that nonsensical or incoherent responses can distort evaluation results~[22,63]. In model merge, parameter interference can further deteriorate this problem by generating invalid responses. If such a response contains no harmful content, it may be classified as safe by automatic evaluators, resulting in an \textit{apparent safety under degeneration}; conversely, depending on the evaluation criteria, the same invalid response may be classified as a successful attack, resulting in an \textit{apparent unsafe under a single evaluator}. Conventional ASR alone therefore cannot distinguish genuine safety improvements from detected generation failures. In fact, our preliminary experiments reveal cases in which the Raw ASR measured by a single evaluator collapses to either 0\% or 100\% under conditions with a high Gibberish Ratio (Appendix~\ref{app:予備実験の詳細結果}).

To address this issue, we introduce a \textit{validity-aware evaluation protocol} that evaluates generated responses along two aspects: whether the response passes rule-based validity checks, and whether it complies with or rejects the harmful request. Specifically, in addition to the conventional Benchmark ASR, we report the Conditional ASR only over valid responses, the Valid Response Rate (VRR), defined as the proportion of valid responses among all generations, and the Valid Safety Rate (VSR), defined as the proportion of responses that are both valid and non-compliant with harmful requests. This protocol enables us to distinguish three cases that conventional ASR-based evaluation conflates: (i) non-degenerate safety refusals, (ii) apparent safety resulting from invalid responses, and (iii) valid responses that nevertheless comply with harmful requests. We further report general utility and perplexity (PPL), allowing us to evaluate merged models from multiple perspectives, including safety, response validity, and task-specific capability.

Re-evaluating conventional merging methods under the Secure Merge setting reveals that some configurations that appear safe under Benchmark ASR alone are instead categorized as exhibiting apparent safety under degeneration or remaining unsafe among non-degenerate responses. Notably, we observe configurations with a Benchmark ASR near 0\% but a VRR near 0\%, as well as configurations with a low Benchmark ASR but a high Conditional ASR. These findings demonstrate that the apparent safety--utility trade-off in Secure Merge, as well as the relative ranking of merging methods, can depend substantially on whether safety is assessed using ASR alone.

To investigate these issues, we address the following research questions (RQs):

\begin{itemize}
    \item \textbf{RQ1 (Validity of Current Evaluation):} To what extent can low benchmark-reported ASR coincide with detected generation failures in the Secure Merge setting?
    
    \item \textbf{RQ2 (Re-evaluating Existing Merging Methods):} How does incorporating VRR and VSR affect the evaluation of representative model merge methods and their safety--utility trade-offs?
    
    \item \textbf{RQ3 (Condition Dependence):} How consistently do validity-aware safety interpretations hold across task domains and merge strengths?
    % How do task specialization, model combinations, intervention granularity, and merging strength influence safety, utility, and response validity under the proposed evaluation metrics?
\end{itemize}

The main contributions of this work are summarized as follows:

\begin{itemize}
    \item We analyze the limitations of ASR-alone evaluation in Secure Merge and show that invalid responses can lead to low benchmark-reported ASR that coincides with severe detected generation failures.
    
    \item We introduce a validity-aware evaluation protocol that combines Benchmark ASR, Conditional ASR, Valid Response Rate (VRR), and Valid Safety Rate (VSR), enabling non-degenerate non-compliant responses, apparent safety caused by invalid responses, and valid but harmful responses to be distinguished.
    
    \item We systematically re-evaluate conventional model merge methods under an identical experimental setting and show that some configurations previously regarded as safe based solely on Benchmark ASR are instead categorized as exhibiting apparent safety under degeneration or remaining unsafe among non-degenerate responses.
    
    \item Through experiments across diverse task domains and merging configurations, we demonstrate that meaningful safety improvements should be assessed jointly using a low Conditional ASR, a high VRR, a high VSR, and preserved downstream utility, rather than ASR alone.
\end{itemize}


% \section{Introduction}
% Large language models (LLMs) achieve strong performance on specialized tasks through fine-tuning (FT), but training and maintaining individual fine-tuned models requires substantial computational resources and data. Model merge has emerged as a cost-effective alternative that integrates multiple fine-tuned models directly in the model parameters, enabling the integration of diverse capabilities of individual models, without additional training [1,2,3]. Since methods such as Model Soups [1] and Task Arithmetic [2,4], a variety of merging methods, including TIES [5], DARE [6], and DELLA [7], have been proposed to mitigate parameter interference among models. However, safety alignment acquired through instruction tuning, RLHF, and related alignment procedures can deteriorate during specialization on benign data and subsequent model merge [8,9,10,11,12]. To address this issue, Secure Merge [13] has been proposed as a framework to restore safety retroactively, in which a safety-specialized model is merged with a task-specialized model. Furthermore, several safety-aware merging methods, including MergeAlign [14], SafeMERGE [15], LED-Merging [16], and Fisher-Weighted Averaging (FWA) [17], have been proposed to mitigate the tradeoffs between safety alignment and task-specific capabilities.

% Secure Merge is evaluated based on the Attack Success Rate (ASR) on harmful requests and utility on downstream tasks [13]. However, we cannot distinguish between valid refusal responses and failures to generate meaningful responses with ASR alone. Parameter interference that occurs during merging may generate invalid responses, such as nonsensical text, repeated special tokens, empty outputs, or regeneration of the input prompt.
% Since ASR does not explicitly assess whether a generated response is a valid and meaningful piece of text, generation failures can reduce the measured ASR. As a result, a model may achieve a low ASR even when its outputs are invalid. Repetitive and degenerate text generation is a well-known manifestation of text degeneration in neural language models~[18].

% The fundamental cause of this issue is that conventional safety evaluators are designed primarily to determine whether a model rejects a harmful request, rather than whether the response itself constitutes a valid and coherent linguistic output [19,20,21]. Prior work on jailbreak evaluation has shown that treating non-refusal alone as evidence of a successful attack does not necessarily align with human judgments and that nonsensical or incoherent responses can distort evaluation results [22,23]. In model merge, parameter interference can further encour1age this problem by generating invalid responses. If such a response contains no harmful content, it may be classified as safe, resulting in a \textit{False Safe}; conversely, depending on the evaluation criteria, the same invalid response may be classified as a successful attack, resulting in a \textit{False Unsafe}. ASR alone therefore cannot distinguish genuine safety improvements from failures in response generation. In fact, our preliminary experiments reveal cases in which the Raw ASR measured by a single evaluator collapses to either 0\% or 100\% under conditions with a high Gibberish Ratio (Appendix~\ref{app:予備実験の詳細結果}).

% To address this issue, we introduce a \textit{validity-aware evaluation metrics} that evaluates generated responses along two aspects: whether the response is valid or invalid, and whether it complies with or rejects the harmful request. Specifically, in addition to the conventional Original ASR, we report the Conditional ASR computed only over valid responses, the Valid Response Rate (VRR), defined as the proportion of valid responses among all generations, and the Valid Safety Rate (VSR), defined as the proportion of responses that are both valid and non-compliant with harmful requests. This protocol enables us to distinguish three cases that conventional ASR-based evaluation conflates: (i) meaningful and valid safety refusals, (ii) apparent safety resulting from invalid responses, and (iii) valid responses that nevertheless comply with harmful requests. We further report downstream utility and perplexity (PPL), allowing us to assess merged models from multiple perspectives, including safety, response validity, and task-specific capability.

% Re-evaluating conventional merging methods under the Secure Merge setting reveals that some configurations that appear safe under Original ASR alone are instead classified as either \textit{False Safe} or \textit{Unsafe but Functional} when response validity is taken into account. Notably, we observe configurations with an Original ASR of 0\% but a VRR of 0\%, as well as configurations with a low Original ASR but a high Conditional ASR. These findings demonstrate that the apparent Safety–Utility trade-off in Secure Merge, as well as the relative ranking of merging methods, can depend substantially on whether safety is assessed using ASR alone.

% To investigate these issues, we address the following research questions (RQs):

% \begin{itemize}
%     \item \textbf{RQ1 (Validity of Current Evaluation):} To what extent does Original ASR misclassify invalid or generation-collapsed responses as safe in the Secure Merge setting?
    
%     \item \textbf{RQ2 (Re-evaluating Existing Merging Methods):} How does incorporating VRR and VSR affect the evaluation of representative model merge methods and their Safety--Utility trade-offs?
    
%     \item \textbf{RQ3 (Condition Dependence):} How do task specialization, model combinations, intervention granularity, and merging strength influence safety, utility, and response validity under the proposed evaluation protocol?
% \end{itemize}

% The main contributions of this work are summarized as follows:

% \begin{itemize}
%     \item We analyze the limitations of ASR-alone evaluation in Secure Merge and show that invalid responses can be misclassified as either \textit{False Safe} or \textit{False Unsafe}, depending on the evaluation criteria.
    
%     \item We introduce a validity-aware evaluation protocol that combines Original ASR, Conditional ASR, Valid Response Rate (VRR), and Valid Safety Rate (VSR), enabling meaningful safety refusals, apparent safety caused by invalid responses, and valid but harmful responses to be distinguished.
    
%     \item We systematically re-evaluate conventional model merge methods under a identical experimental setting and show that some configurations previously regarded as safe based solely on Original ASR are instead categorized as \textit{False Safe} or \textit{Unsafe but Functional}.
    
%     \item Through experiments across diverse task domains and merging configurations, we demonstrate that meaningful safety improvements be assessed jointly using a low Conditional ASR, a high VRR, a high VSR, and preserved downstream utility, rather than ASR alone.
% \end{itemize}

\section{Related Work}

\subsection{Model Merge and Interference Mitigation}

Model Soups~[1] and Task Arithmetic~[2,4] are foundational model merge methods that combine either the weights of multiple fine-tuned models or their weight differences with the base model, which realizes multiple task capabilities without additional training. However, simply merging multiple models often causes parameter interference, where conflicting parameter updates cancel or overwrite one another, leading to degradation of the specialized capabilities of the merged model. To mitigate this problem, TIES~[5] resolves sign conflicts among task vectors, while DARE~[6] and DELLA~[7] reduce interference through parameter sparsification and importance-aware parameter selection. These methods primarily focus on sustaining task performance across multiple domains rather than maintaining safety alignment.

\subsection{Secure Merge and Safety-Preserving Model Merge}

To address the degradation of safety alignment caused by fine-tuning and model merge, Secure Merge~[13] introduces a framework to restore safety retroactively by merging a safety-specialized model, referred to as a \emph{safety patch}, into the task-specialized model. While this approach effectively restores safety, interference between the safety patch and the task-specific model may arise when both share the same parameters, giving rise to the well-known \emph{Safety Tax}, in which improved safety comes at the expense of downstream utility~[13,24]. Hammoud et al.~[8] further reported the ``one bad model spoils the bunch'' phenomenon, demonstrating that introducing a single misaligned expert model can substantially compromise the safety of the merged model and that general interference-mitigation methods such as TIES are insufficient to prevent this failure.

To mitigate the interference between safety and utility, a variety of merging methods have been proposed. MergeAlign~[14] tunes task vectors associated with domain-specific capabilities and safety alignment, whereas SafeMERGE~[15] and LED-Merging~[16] resolve safety--utility trade-offs at the layer and neuron levels, respectively. Fisher-Weighted Averaging (FWA)~[17] estimates parameter importance using the diagonal Fisher information matrix and performs importance-weighted model merging. AlignMerge~[25] further leverages Fisher geometry and alignment subspaces to keep parameter space that are vital for maintaining safety.
Fisher-based merging has also been extended in several directions, including uncertainty-aware gradient matching, Fisher mask nodes, Fisher-weighted median merging, Bayesian optimization for adaptive Fisher weighting, and data-free layer-adaptive Fisher merging~[26,27,28,29,30].
Despite their methodological differences, these approaches primarily focus on designing merging rules or intervention strategies that mitigate interference between safety and utility. Their effectiveness is typically evaluated in terms of reductions in Attack Success Rate (ASR) together with utility on downstream tasks~[8,13,14,15,16].


\subsection{Safety Evaluation for Large Language Models}
The Attack Success Rate (ASR), which evaluates a model's behavior with harmful requests under adversarial prompts, is the standard metric for evaluating the jailbreak robustness of large language models. HarmBench~[19] provides an evaluation pipeline that defines ASR by determining whether a model's output successfully achieves the intended harmful compliance. Its classifier is designed primarily to assess harmful-behavior completion rather than to evaluate response validity as an independent dimension or systematically identify generation failures, such as empty outputs, gibberish, repetitive special tokens, or verbatim prompt repetition. Similarly, Llama Guard~[20] is an LLM-based safety classifier that categorizes model inputs and outputs as \textit{safe} or \textit{unsafe} according to predefined safety policies.


Several recent studies have sought to address limitations of binary compliance-based jailbreak evaluation. StrongREJECT evaluates refusal behavior beyond simple non-refusal by considering the usefulness and strength of a model's response to harmful requests~[22]. WildGuard jointly evaluates harmfulness, jailbreak behavior, and refusal behavior within a unified moderation framework~[21]. More recently, JailMeter proposes an evidence-based evaluation framework in which an attack is considered successful only when the response captures the malicious intent and provides a complete answer, thereby reducing false-positive attack-success judgments caused by incomplete or ambiguous outputs~[63].


These approaches primarily aim to improve the reliability or granularity of harmful-compliance and refusal judgments. In contrast, we focus on the Secure Merge setting, where parameter interference can cause severe generation failures. Rather than proposing a new harmfulness evaluator, we treat response validity as an independent evaluation axis and quantify detected severe degeneration through rule-based checks, Valid Response Rate (VRR), Conditional ASR, and Valid Safety Rate (VSR). Additional background and discussion are provided in Appendix~\ref{app:本研究の背景と定式化上の位置づけ}.


\section{Evaluation}

This section describes the experimental setup used throughout this paper. Following the Secure Merge setting, we operationalize \emph{safety} as robustness against jailbreak attacks and evaluate how effectively each merging method preserves this property while maintaining downstream task utility.


\subsection{Models and Baselines}

We use \texttt{Llama-2-7B}~[31] as the base model and consider three task-specialized models covering mathematics, code generation, and the medical domain~[32,33,34], together with a safety-specialized model fine-tuned for jailbreak robustness~[35]. We compare representative model merge methods spanning standard merging (Task Arithmetic), interference-aware merging (TIES, DARE, and DELLA), safety-preserving merging (MergeAlign, SafeMERGE, and LED-Merging), and Fisher importance-based merging (FWA).

\begin{itemize}
    \item \textbf{Base model ($\theta_0$):} Pretrained \texttt{Llama-2-7B} checkpoint~[31] (exact identifier: \texttt{meta-llama/Llama-2-7b-hf}).

    \item \textbf{Task-specialized models ($\theta_u$):} We use WizardMath-7B-V1.0~[32], WizardCoder-Python-7B-V1.0~[33], and MedAlpaca-7B~[34] for the mathematics, code generation, and medical domains, respectively. Each model follows the original training configuration described in its model card.

    \item \textbf{Calibration data for utility Fisher estimation:} We use calibration samples from GSM8K for mathematics, Evol-Instruct-Code for code generation, and MedAlpaca for the medical domain.


    \item \textbf{Safety-specialized model ($\theta_s$):} We fine-tune pretrained \texttt{Llama-2-7B} using LoRA (rank $r=16$, $\alpha=32$) to obtain a safety-aligned model. The fine-tuning dataset consists of 1,400 harmful prompts from the \textit{JAILBREAK TRIGGER} subset of TrustLLM~[35], where each prompt is paired with the exact same fixed refusal response: \texttt{"I'm sorry, but I cannot assist with that request."} Standard causal language modeling (CLM) loss is computed across all tokens, including both prompt and completion tokens. This controlled setup reduces variation arising from refusal phrasing. However, the fixed refusal target and prompt-inclusive loss may increase susceptibility to verbatim prompt repetition or formulaic template copying under severe parameter interference.
    
    (discussed in Appendix~\ref{app:実験設計の詳細}).
\end{itemize}

The objective of this work is not to establish the superiority of any particular merging method, but to provide a consistent evaluation of representative approaches with different intervention granularities and information sources under a common experimental setting. We evaluate MergeAlign~[14] as a representative safety-aware baseline. In contrast, we do not evaluate AlignMerge~[25] since an official implementation was not publicly available at the time of our study. Detailed formulations and implementation settings for all merging methods are provided in Appendix~\ref{app:実験設計の詳細}.

\subsection{Evaluation Protocol}

For safety evaluation, we use four benchmark datasets: HarmBench~[19], JailbreakBench~[36], StrongREJECT~[22], and WildJailbreak~[37]. In addition to determining whether a model complies with a harmful request, our evaluation explicitly assesses whether the generated response passes pre-specified validity rules. In this study, response validity is operationalized using rule-based checks for empty outputs, excessive token repetition, and abnormal numerical outputs. Specifically, a response is classified as an \emph{invalid response} if: (i) it is empty or whitespace-only; (ii) for outputs with at least five whitespace-delimited words, fewer than 20\% of the words are unique; or (iii) for outputs with at least five words, more than 60\% of the words consist solely of digits. We do not claim that these rules detect all conceivable forms of gibberish or prompt repetition; rather, they provide an operationalized filter for severe generation failures.

Based on the above protocol, we distinguish between standard benchmark-reported evaluations and our proposed validity-aware metrics:

\begin{itemize}
    \item \textbf{Benchmark ASR ($\mathrm{ASR}_{\mathrm{bench}}$)}: The standard Attack Success Rate computed directly using each benchmark's official evaluator pipeline over all $N$ generated responses: $\mathrm{ASR}_{\mathrm{bench}} = \frac{1}{N}\sum_{i=1}^N B_i$, where $B_i \in \{0,1\}$ is the binary attack-success decision assigned by the benchmark classifier.
    
    \item \textbf{Validity-aware Metrics}: For each response sample $i \in \{1,\dots,N\}$, let $V_i \in \{0,1\}$ be the binary indicator of whether the output passes our rule-based validity checks, and let $H_i \in \{0,1\}$ be the harmfulness decision assigned by a common evaluator (HarmBench Classifier) over valid outputs. We define:
    \begin{itemize}
        \item \textbf{Conditional ASR ($\mathrm{ASR}_{\mathrm{valid}}$)}: The Attack Success Rate restricted to non-degenerate responses ($N_{\mathrm{valid}} = \sum V_i$ samples): $\mathrm{ASR}_{\mathrm{valid}} = \frac{\sum_{i=1}^N V_i H_i}{\sum_{i=1}^N V_i}$.
        \item \textbf{Valid Response Rate (VRR)}: The proportion of valid responses among all generated outputs: $\mathrm{VRR} = \frac{1}{N}\sum_{i=1}^N V_i$.
        \item \textbf{Valid Safety Rate (VSR)}: The proportion of responses that are both valid and non-compliant with harmful requests: $\mathrm{VSR} = \frac{1}{N}\sum_{i=1}^N V_i(1-H_i)$.
    \end{itemize}
    \item \textbf{Utility}: Downstream task performance measured on GSM8K~[38], Math500~[39], HumanEval~[40], MBPP~[41], MedQA~[42], PubMedQA~[43], AlpacaEval~[44], IFEval~[45], and MMLU~[46]. For each domain, utility is reported as the unweighted arithmetic mean of the percentage scores of its constituent benchmarks. Detailed per-benchmark scores are provided in Appendix~\ref{app:主実験の詳細結果}.
\end{itemize}

Under these validity-aware metrics, for any individual evaluation run, the metrics satisfy the exact mathematical identity:
\[
\mathrm{VSR} = \mathrm{VRR}\left(1-\mathrm{ASR}_{\mathrm{valid}}\right)
\]
% Note that since $\mathrm{ASR}_{\mathrm{bench}}$ is derived from official benchmark evaluators whereas $H_i$ is evaluated on valid outputs, $\mathrm{ASR}_{\mathrm{bench}}$ and $\mathrm{ASR}_{\mathrm{valid}}$ are evaluated separately, allowing us to reveal cases where a low $\mathrm{ASR}_{\mathrm{bench}}$ coincides with generation failures rather than true safety preservation.


Because $\mathrm{ASR}_{\mathrm{bench}}$ is obtained using each benchmark's official evaluator over all outputs, whereas $\mathrm{ASR}_{\mathrm{valid}}$ uses a common HarmBench Classifier only over outputs that pass the rule-based validity checks, the two quantities are not numerically interchangeable. Their discrepancy should be interpreted diagnostically: it can reflect both conditioning on detected valid outputs and differences between evaluator pipelines. We therefore use these metrics jointly to identify cases in which low benchmark-reported ASR co-occurs with low VRR or high harmful compliance among outputs passing the validity checks.

For each benchmark, the generation prompt template, official benchmark evaluator, and harmful-compliance mapping rules are summarized in Appendix~\ref{app:実験設計の詳細}.

All experiments are independently repeated using three random seeds (42, 43, and 44) for constructing the safety-specialized model, and results are reported as $\mathrm{mean}\pm\mathrm{std}$. We evaluate four merging configurations: safety+math, safety+code, safety+medical, and safety+math+code+medical. Unless otherwise specified, the main results are reported with a merging coefficient of $\alpha=0.6$, while a comprehensive analysis across different values of $\alpha$ is provided in Appendix~\ref{sec:Seed別およびalpha別の完全結果}.



\input{Formatting_Instructions_For_NeurIPS_2026/tables/table1}

\subsection{From Preliminary Experiments to the Main Evaluation}

As a preliminary investigation, we conducted an exploratory evaluation using TrustLLM~[35] and BeaverTails~[47] with a single safety evaluator. Table~\ref{tab:prelim_detail} presents the results for the \texttt{safety+math} configuration (see Appendix~\ref{app:予備実験の詳細結果} for the complete results). Under conditions with a high Gibberish Ratio, the Raw ASR measured by TrustLLM frequently collapsed to either 0\% or 100\%. This observation suggests that generation failures can substantially affect the ASR reported by a single evaluator, although the preliminary experiment does not establish whether individual evaluator judgments are erroneous.

Motivated by these findings, the main experiments do not rely on ASR alone. Instead, we explicitly separate harmfulness assessment from response validity assessment. This evaluation protocol enables us to distinguish whether a reduction in ASR reflects valid non-compliant responses or merely results from generation failures.





\section{Results and Discussion}

This section presents the main findings from our validity-aware evaluation. Detailed experimental results and supplementary analyses are provided in Appendix~\ref{app:予備実験の詳細結果}, \ref{app:主実験の詳細結果}, \ref{app:偽の安全性の追加分析}, \ref{app:計算資源とデータ要件}, and \ref{sec:Seed別およびalpha別の完全結果}.

Table~\ref{tab:safety_diagnostic} presents representative examples in which the interpretation of safety changes after applying the proposed validity-aware evaluation protocol. The table reports Benchmark ASR ($\mathrm{ASR}_{\mathrm{bench}}$), Conditional ASR ($\mathrm{ASR}_{\mathrm{valid}}$), Valid Response Rate (VRR), Valid Safety Rate (VSR), and utility. These examples show that conditions exhibiting similarly low Benchmark ASR can correspond to substantially different safety states once response validity is taken into account.

MergeAlign achieves a Benchmark ASR of 0.00\% under multiple settings. However, both its VRR and VSR are also 0.00\%, indicating that it fails to generate any responses that pass our validity checks. Based on Benchmark ASR alone, this result suggests that the model successfully rejects all harmful requests; in practice, however, none of the evaluated responses are valid. In this case, the apparent reduction in ASR reflects a loss of generation capability rather than improved safety. We therefore categorize this behavior as \textit{Apparent safety under degeneration}, where generation failure or invalid responses coincide with a low benchmark ASR. This result demonstrates that a low ASR, while potentially necessary for establishing safety improvement, is not sufficient on its own.

\input{Formatting_Instructions_For_NeurIPS_2026/tables/table3}

Benchmark-reported ASR can diverge substantially from harmful compliance measured among outputs passing our validity checks.  Under the \texttt{safety+code} setting, methods such as \texttt{matena\_fisher} achieve a Benchmark ASR of only 3.37\%, whereas its Conditional ASR reaches 65.04\% with a VRR of 59.73\%. Thus, among non-degenerate responses, harmful compliance is dramatically higher than implied by the raw unconditioned benchmark ASR.

In contrast, Task Arithmetic under the \texttt{safety+math} setting exhibits low Benchmark ASR (2.48\%) and Conditional ASR (3.88\%) together with high VRR (99.92\%) and VSR (96.04\%). In this case, the low Benchmark ASR co-occurs with high VRR and low Conditional ASR, making this configuration consistent with non-degenerate safety behavior under our evaluation protocol rather than with the detected generation failures. Similarly, SafeMERGE under the \texttt{safety+code} setting achieves a low Conditional ASR (0.34\%) together with high VRR (99.92\%) and VSR (99.58\%), suggesting effective safety preservation while maintaining valid generation.

Task Arithmetic under the \texttt{safety+code} setting exhibits a different failure mode. Its Benchmark ASR is 29.17\%, whereas its Conditional ASR is 60.91\%, with a VSR of 36.94\%. Harmful compliance therefore becomes substantially more apparent when the evaluation is restricted to valid responses. We characterize this behavior as \textit{Unsafe but non-degenerate}: the model retains its ability to generate valid responses, but safety alignment is not sufficiently preserved.

Taken together, these results show that Benchmark ASR alone cannot distinguish among at least three qualitatively different outcomes: (i) non-degenerate and safe behavior, (ii) apparent safety under detected generation failures, and (iii) preserved generation capability with substantial residual harmful compliance. We therefore argue that safety evaluation in Secure Merge should report Conditional ASR, VRR, and VSR alongside Benchmark ASR to distinguish genuine safety preservation from artifacts of generation failure.





\subsection{Safety--Utility Trade-offs Under Validity-Aware Evaluation}


To generalize these observations, we categorize all methods and experimental settings into four groups: \textit{Non-degenerate and safe}, \textit{Apparent safety under degeneration}, \textit{Apparent unsafe under a single evaluator}, and \textit{Unsafe but non-degenerate} (Table~\ref{tab:safety_diagnostic}). 
% A \textit{Non-degenerate and safe} configuration simultaneously achieves a low Conditional ASR, high VRR, high VSR, and preserved downstream utility. 
A \textit{Non-degenerate and safe} configuration exhibits low Conditional
ASR together with high VRR and VSR. We report downstream utility separately,
because the degree of utility retention depends on the reference model,
task domain, and deployment objective. An \textit{Apparent safety under degeneration} configuration exhibits a low Benchmark ASR but poor response validity, such as a low VRR or abnormal PPL, indicating that the apparent safety originates from generation failures. An \textit{Apparent unsafe under a single evaluator} configuration refers to cases in which a large number of invalid responses are judged as successful attacks by a single evaluator; in our experiments, this behavior is primarily observed in the preliminary evaluation using a single evaluator. Finally, an \textit{Unsafe but non-degenerate} configuration maintains relatively high response validity but also exhibits a high Conditional ASR, indicating that the model continues to produce valid responses while remaining vulnerable to harmful requests.


This validity-aware categorization substantially changes the interpretation of the Safety--Utility trade-off in Secure Merge. Under the conventional ASR--Utility perspective, configurations suffering from generation collapse may appear favorable simply because their ASR is artificially reduced. By incorporating VRR and VSR, these cases are correctly identified as failures of response generation rather than successful safety preservation. Response validity therefore serves not only as a third evaluation dimension alongside safety and utility, but also as a prerequisite for interpreting low ASR as evidence of genuine safety improvement.

More broadly, the proposed validity-aware evaluation enables meaningful comparisons among model merge methods by considering not only safety but also whether valid generation is preserved. This perspective is particularly important in practice when selecting merging strategies, intervention strengths, and safety patches for deployment.


\subsection{Limitations and Scope}

One limitation of this work is that the distinction between valid and invalid responses depends on the predefined validity criteria. In our evaluation, severe generation failures are operationalized as (i) empty or whitespace-only outputs; (ii) for outputs with at least five whitespace-delimited words, fewer than 20\% unique words; or (iii) for outputs with at least five words, more than 60\% numeric-only words. These criteria are intended to identify severe generation failures rather than to exhaustively characterize all forms as gibberish, special-token repetition, or prompt repetition. Accordingly, the reported VRR and VSR should be interpreted with respect to this operational definition of response validity.

A second limitation concerns the automatic harmfulness evaluators and the rule-based response-validity checks used in our evaluation. Misclassification may still occur when automatic classifiers fail to properly assess partially compliant or ambiguous responses. Future work should validate the proposed validity criteria through human annotation, analyze inter-evaluator agreement, and consider introducing an \textit{undetermined} category for responses that cannot be reliably classified.

Finally, our experiments focus on Llama-2-7B and do not evaluate the generality of the proposed evaluation metrics across larger language models, alternative model architectures, or parameter-efficient adaptation methods such as LoRA. In addition, due to computational constraints, each benchmark is evaluated using at most 320 samples. Accordingly, this work does not seek to establish a definitive ranking of existing merging methods or to claim generalization across model scales. Instead, its primary contribution is to demonstrate the importance of explicitly evaluating response validity in the Secure Merge setting.

\section{Conclusion}

In this work, we showed that low benchmark-reported Attack Success Rate (ASR) does not necessarily indicate successful safety preservation in Secure Merge, making ASR alone insufficient to determine whether a merged model remains both safe and capable of producing non-degenerate responses. Therefore, we introduced a validity-aware evaluation protocol that jointly reports harmful compliance, rule-based response validity, and general utility.

Our experiments reveal qualitatively distinct safety regimes that can appear similar under ASR alone, indicating genuinely safe behavior, apparent safety caused by generation failures, and functional responses that still exhibit substantial harmful compliance. These results suggest that safety evaluation for Secure Merge should report validity-conditioned harmful compliance and response-validity statistics alongside conventional benchmark ASR.

Our primary contribution is an evaluation perspective rather than a new merging algorithm. Future work should validate the proposed rule-based validity criteria through human annotation, analyze agreement among safety evaluators, and extend the evaluation to larger model families and additional adaptation settings.




\section*{References}
\small
\input{Formatting_Instructions_For_NeurIPS_2026/bib}

\newpage
\appendix

\input{Formatting_Instructions_For_NeurIPS_2026/appendix}

\clearpage


\newpage

\section*{NeurIPS Paper Checklist}
%%% BEGIN INSTRUCTIONS %%%
% The checklist is designed to encourage best practices for responsible machine learning research, addressing issues of reproducibility, transparency, research ethics, and societal impact. Do not remove the checklist: {\bf The papers not including the checklist will be desk rejected.} The checklist should follow the references and follow the (optional) supplemental material.  The checklist does NOT count towards the page
% limit. 

% Please read the checklist guidelines carefully for information on how to answer these questions. For each question in the checklist:
% \begin{itemize}
%     \item You should answer \answerYes{}, \answerNo{}, or \answerNA{}.
%     \item \answerNA{} means either that the question is Not Applicable for that particular paper or the relevant information is Not Available.
%     \item Please provide a short (1--2 sentence) justification right after your answer (even for \answerNA). 
%    % \item {\bf The papers not including the checklist will be desk rejected.}
% \end{itemize}

% {\bf The checklist answers are an integral part of your paper submission.} They are visible to the reviewers, area chairs, senior area chairs, and ethics reviewers. You will also be asked to include it (after eventual revisions) with the final version of your paper, and its final version will be published with the paper.

% The reviewers of your paper will be asked to use the checklist as one of the factors in their evaluation. While \answerYes{} is generally preferable to \answerNo{}, it is perfectly acceptable to answer \answerNo{} provided a proper justification is given (e.g., error bars are not reported because it would be too computationally expensive'' or ``we were unable to find the license for the dataset we used''). In general, answering \answerNo{} or \answerNA{} is not grounds for rejection. While the questions are phrased in a binary way, we acknowledge that the true answer is often more nuanced, so please just use your best judgment and write a justification to elaborate. All supporting evidence can appear either in the main paper or the supplemental material, provided in appendix. If you answer \answerYes{} to a question, in the justification please point to the section(s) where related material for the question can be found.

% IMPORTANT, please:
% \begin{itemize}
%     \item {\bf Delete this instruction block, but keep the section heading ``NeurIPS Paper Checklist"},
%     \item  {\bf Keep the checklist subsection headings, questions/answers and guidelines below.}
%     \item {\bf Do not modify the questions and only use the provided macros for your answers}.
% \end{itemize} 
 

% %%% END INSTRUCTIONS %%%

\begin{enumerate}

\item {\bf Claims}
    \item[] Question: Do the main claims made in the abstract and introduction accurately reflect the paper's contributions and scope?
    \item[] Answer: \answerYes{} % Replace by \answerYes{}, \answerNo{}, or \answerNA{}.
    \item[] Justification:  The abstract and Section 1 state the paper's scope as a validity-aware evaluation of Secure Merge, rather than as a new merging algorithm. The claims are supported by the experimental results in Sections 3 and 4
    \item[] Guidelines:
    \begin{itemize}
        \item The answer \answerNA{} means that the abstract and introduction do not include the claims made in the paper.
        \item The abstract and/or introduction should clearly state the claims made, including the contributions made in the paper and important assumptions and limitations. A \answerNo{} or \answerNA{} answer to this question will not be perceived well by the reviewers. 
        \item The claims made should match theoretical and experimental results, and reflect how much the results can be expected to generalize to other settings. 
        \item It is fine to include aspirational goals as motivation as long as it is clear that these goals are not attained by the paper. 
    \end{itemize}

\item {\bf Limitations}
    \item[] Question: Does the paper discuss the limitations of the work performed by the authors?
    \item[] Answer: \answerYes{} % Replace by \answerYes{}, \answerNo{}, or \answerNA{}.
    \item[] Justification: Section 4.2 and Appendix C.3 discuss limitations of the rule-based validity definition, automatic harmfulness evaluation, possible training--evaluation overlap, restricted model scale, benchmark subsampling, and the fixed-refusal safety-training design.
    \item[] Guidelines:
    \begin{itemize}
        \item The answer \answerNA{} means that the paper has no limitation while the answer \answerNo{} means that the paper has limitations, but those are not discussed in the paper. 
        \item The authors are encouraged to create a separate ``Limitations'' section in their paper.
        \item The paper should point out any strong assumptions and how robust the results are to violations of these assumptions (e.g., independence assumptions, noiseless settings, model well-specification, asymptotic approximations only holding locally). The authors should reflect on how these assumptions might be violated in practice and what the implications would be.
        \item The authors should reflect on the scope of the claims made, e.g., if the approach was only tested on a few datasets or with a few runs. In general, empirical results often depend on implicit assumptions, which should be articulated.
        \item The authors should reflect on the factors that influence the performance of the approach. For example, a facial recognition algorithm may perform poorly when image resolution is low or images are taken in low lighting. Or a speech-to-text system might not be used reliably to provide closed captions for online lectures because it fails to handle technical jargon.
        \item The authors should discuss the computational efficiency of the proposed algorithms and how they scale with dataset size.
        \item If applicable, the authors should discuss possible limitations of their approach to address problems of privacy and fairness.
        \item While the authors might fear that complete honesty about limitations might be used by reviewers as grounds for rejection, a worse outcome might be that reviewers discover limitations that aren't acknowledged in the paper. The authors should use their best judgment and recognize that individual actions in favor of transparency play an important role in developing norms that preserve the integrity of the community. Reviewers will be specifically instructed to not penalize honesty concerning limitations.
    \end{itemize}

\item {\bf Theory assumptions and proofs}
    \item[] Question: For each theoretical result, does the paper provide the full set of assumptions and a complete (and correct) proof?
    \item[] Answer: \answerNA{} % Replace by \answerYes{}, \answerNo{}, or \answerNA{}.
    \item[] Justification: This paper does not present theoretical theorems or formal proofs. Its contributions are empirical and methodological.
    \item[] Guidelines:
    \begin{itemize}
        \item The answer \answerNA{} means that the paper does not include theoretical results. 
        \item All the theorems, formulas, and proofs in the paper should be numbered and cross-referenced.
        \item All assumptions should be clearly stated or referenced in the statement of any theorems.
        \item The proofs can either appear in the main paper or the supplemental material, but if they appear in the supplemental material, the authors are encouraged to provide a short proof sketch to provide intuition. 
        \item Inversely, any informal proof provided in the core of the paper should be complemented by formal proofs provided in appendix or supplemental material.
        \item Theorems and Lemmas that the proof relies upon should be properly referenced. 
    \end{itemize}

    \item {\bf Experimental result reproducibility}
    \item[] Question: Does the paper fully disclose all the information needed to reproduce the main experimental results of the paper to the extent that it affects the main claims and/or conclusions of the paper (regardless of whether the code and data are provided or not)?
    \item[] Answer: \answerYes{} % Replace by \answerYes{}, \answerNo{}, or \answerNA{}.
    \item[] Justification: Sections 3 and Appendix C specify the model checkpoints, safety-data source, LoRA configuration, optimizer, training schedule, random seeds, merge methods, calibration data, benchmark subsets, decoding settings, validity rules, and aggregation procedure used for the reported results.
    \item[] Guidelines:
    \begin{itemize}
        \item The answer \answerNA{} means that the paper does not include experiments.
        \item If the paper includes experiments, a \answerNo{} answer to this question will not be perceived well by the reviewers: Making the paper reproducible is important, regardless of whether the code and data are provided or not.
        \item If the contribution is a dataset and\slash or model, the authors should describe the steps taken to make their results reproducible or verifiable. 
        \item Depending on the contribution, reproducibility can be accomplished in various ways. For example, if the contribution is a novel architecture, describing the architecture fully might suffice, or if the contribution is a specific model and empirical evaluation, it may be necessary to either make it possible for others to replicate the model with the same dataset, or provide access to the model. In general. releasing code and data is often one good way to accomplish this, but reproducibility can also be provided via detailed instructions for how to replicate the results, access to a hosted model (e.g., in the case of a large language model), releasing of a model checkpoint, or other means that are appropriate to the research performed.
        \item While NeurIPS does not require releasing code, the conference does require all submissions to provide some reasonable avenue for reproducibility, which may depend on the nature of the contribution. For example
        \begin{enumerate}
            \item If the contribution is primarily a new algorithm, the paper should make it clear how to reproduce that algorithm.
            \item If the contribution is primarily a new model architecture, the paper should describe the architecture clearly and fully.
            \item If the contribution is a new model (e.g., a large language model), then there should either be a way to access this model for reproducing the results or a way to reproduce the model (e.g., with an open-source dataset or instructions for how to construct the dataset).
            \item We recognize that reproducibility may be tricky in some cases, in which case authors are welcome to describe the particular way they provide for reproducibility. In the case of closed-source models, it may be that access to the model is limited in some way (e.g., to registered users), but it should be possible for other researchers to have some path to reproducing or verifying the results.
        \end{enumerate}
    \end{itemize}


\item {\bf Open access to data and code}
    \item[] Question: Does the paper provide open access to the data and code, with sufficient instructions to faithfully reproduce the main experimental results, as described in supplemental material?
    \item[] Answer: \answerNo{} % Replace by \answerYes{}, \answerNo{}, or \answerNA{}.
    \item[] Justification: We do not provide code, trained safety-model checkpoints, or generated outputs with this submission. The paper documents the experimental protocol and uses publicly available base models, task-specialized models, benchmarks, and merging implementations.
    \item[] Guidelines:
    \begin{itemize}
        \item The answer \answerNA{} means that paper does not include experiments requiring code.
        \item Please see the NeurIPS code and data submission guidelines (\url{https://neurips.cc/public/guides/CodeSubmissionPolicy}) for more details.
        \item While we encourage the release of code and data, we understand that this might not be possible, so \answerNo{} is an acceptable answer. Papers cannot be rejected simply for not including code, unless this is central to the contribution (e.g., for a new open-source benchmark).
        \item The instructions should contain the exact command and environment needed to run to reproduce the results. See the NeurIPS code and data submission guidelines (\url{https://neurips.cc/public/guides/CodeSubmissionPolicy}) for more details.
        \item The authors should provide instructions on data access and preparation, including how to access the raw data, preprocessed data, intermediate data, and generated data, etc.
        \item The authors should provide scripts to reproduce all experimental results for the new proposed method and baselines. If only a subset of experiments are reproducible, they should state which ones are omitted from the script and why.
        \item At submission time, to preserve anonymity, the authors should release anonymized versions (if applicable).
        \item Providing as much information as possible in supplemental material (appended to the paper) is recommended, but including URLs to data and code is permitted.
    \end{itemize}


\item {\bf Experimental setting/details}
    \item[] Question: Does the paper specify all the training and test details (e.g., data splits, hyperparameters, how they were chosen, type of optimizer) necessary to understand the results?
    \item[] Answer: \answerYes{} % Replace by \answerYes{}, \answerNo{}, or \answerNA{}.
    \item[] Justification: The paper specifies all training and evaluation details in Section 3 and Appendix C.
    \item[] Guidelines:
    \begin{itemize}
        \item The answer \answerNA{} means that the paper does not include experiments.
        \item The experimental setting should be presented in the core of the paper to a level of detail that is necessary to appreciate the results and make sense of them.
        \item The full details can be provided either with the code, in appendix, or as supplemental material.
    \end{itemize}

\item {\bf Experiment statistical significance}
    \item[] Question: Does the paper report error bars suitably and correctly defined or other appropriate information about the statistical significance of the experiments?
    \item[] Answer: \answerYes{} % Replace by \answerYes{}, \answerNo{}, or \answerNA{}.
    \item[] Justification: Error bars are suitably and correctly defined in the paper.
    \item[] Guidelines:
    \begin{itemize}
        \item The answer \answerNA{} means that the paper does not include experiments.
        \item The authors should answer \answerYes{} if the results are accompanied by error bars, confidence intervals, or statistical significance tests, at least for the experiments that support the main claims of the paper.
        \item The factors of variability that the error bars are capturing should be clearly stated (for example, train/test split, initialization, random drawing of some parameter, or overall run with given experimental conditions).
        \item The method for calculating the error bars should be explained (closed form formula, call to a library function, bootstrap, etc.)
        \item The assumptions made should be given (e.g., Normally distributed errors).
        \item It should be clear whether the error bar is the standard deviation or the standard error of the mean.
        \item It is OK to report 1-sigma error bars, but one should state it. The authors should preferably report a 2-sigma error bar than state that they have a 96\% CI, if the hypothesis of Normality of errors is not verified.
        \item For asymmetric distributions, the authors should be careful not to show in tables or figures symmetric error bars that would yield results that are out of range (e.g., negative error rates).
        \item If error bars are reported in tables or plots, the authors should explain in the text how they were calculated and reference the corresponding figures or tables in the text.
    \end{itemize}

\item {\bf Experiments compute resources}
    \item[] Question: For each experiment, does the paper provide sufficient information on the computer resources (type of compute workers, memory, time of execution) needed to reproduce the experiments?
    \item[] Answer: \answerNo{} % Replace by \answerYes{}, \answerNo{}, or \answerNA{}.
    \item[] Justification: Appendix H reports merge-time and peak-memory measurements for conventional methods on an NVIDIA h100 GPU, but we do not provide complete training and end-to-end evaluation compute for every reported experiment.
    \item[] Guidelines:
    \begin{itemize}
        \item The answer \answerNA{} means that the paper does not include experiments.
        \item The paper should indicate the type of compute workers CPU or GPU, internal cluster, or cloud provider, including relevant memory and storage.
        \item The paper should provide the amount of compute required for each of the individual experimental runs as well as estimate the total compute. 
        \item The paper should disclose whether the full research project required more compute than the experiments reported in the paper (e.g., preliminary or failed experiments that didn't make it into the paper). 
    \end{itemize}
    
\item {\bf Code of ethics}
    \item[] Question: Does the research conducted in the paper conform, in every respect, with the NeurIPS Code of Ethics \url{https://neurips.cc/public/EthicsGuidelines}?
    \item[] Answer: \answerYes{} % Replace by \answerYes{}, \answerNo{}, or \answerNA{}.
    \item[] Justification: We have checked the code of ethics, and the research is confirmed in every respect.
    \item[] Guidelines:
    \begin{itemize}
        \item The answer \answerNA{} means that the authors have not reviewed the NeurIPS Code of Ethics.
        \item If the authors answer \answerNo, they should explain the special circumstances that require a deviation from the Code of Ethics.
        \item The authors should make sure to preserve anonymity (e.g., if there is a special consideration due to laws or regulations in their jurisdiction).
    \end{itemize}


\item {\bf Broader impacts}
    \item[] Question: Does the paper discuss both potential positive societal impacts and negative societal impacts of the work performed?
    \item[] Answer: \answerYes{} % Replace by \answerYes{}, \answerNo{}, or \answerNA{}.
    \item[] Justification: This work can reduce the risk that non-functional merged models are mistakenly judged as safe and deployed without adequate validation. However, the evaluation protocol may also be misused to optimize for specific benchmark criteria or conceal failures not captured by the rule-based validity checks. We therefore recommend reporting the full set of safety, validity, and utility metrics, publishing the exact validity rules, and complementing automatic evaluation with human
audits for deployment-relevant settings.
    \item[] Guidelines:
    \begin{itemize}
        \item The answer \answerNA{} means that there is no societal impact of the work performed.
        \item If the authors answer \answerNA{} or \answerNo, they should explain why their work has no societal impact or why the paper does not address societal impact.
        \item Examples of negative societal impacts include potential malicious or unintended uses (e.g., disinformation, generating fake profiles, surveillance), fairness considerations (e.g., deployment of technologies that could make decisions that unfairly impact specific groups), privacy considerations, and security considerations.
        \item The conference expects that many papers will be foundational research and not tied to particular applications, let alone deployments. However, if there is a direct path to any negative applications, the authors should point it out. For example, it is legitimate to point out that an improvement in the quality of generative models could be used to generate Deepfakes for disinformation. On the other hand, it is not needed to point out that a generic algorithm for optimizing neural networks could enable people to train models that generate Deepfakes faster.
        \item The authors should consider possible harms that could arise when the technology is being used as intended and functioning correctly, harms that could arise when the technology is being used as intended but gives incorrect results, and harms following from (intentional or unintentional) misuse of the technology.
        \item If there are negative societal impacts, the authors could also discuss possible mitigation strategies (e.g., gated release of models, providing defenses in addition to attacks, mechanisms for monitoring misuse, mechanisms to monitor how a system learns from feedback over time, improving the efficiency and accessibility of ML).
    \end{itemize}
    
\item {\bf Safeguards}
    \item[] Question: Does the paper describe safeguards that have been put in place for responsible release of data or models that have a high risk for misuse (e.g., pre-trained language models, image generators, or scraped datasets)?
    \item[] Answer: \answerNA{} % Replace by \answerYes{}, \answerNo{}, or \answerNA{}.
    \item[] Justification: This work does not introduce or release a new generative model, harmful-completion dataset, or capability that materially increases misuse risk. We evaluate existing models using established safety benchmarks and report only limited illustrative excerpts. 
    \item[] Guidelines:
    \begin{itemize}
        \item The answer \answerNA{} means that the paper poses no such risks.
        \item Released models that have a high risk for misuse or dual-use should be released with necessary safeguards to allow for controlled use of the model, for example by requiring that users adhere to usage guidelines or restrictions to access the model or implementing safety filters. 
        \item Datasets that have been scraped from the Internet could pose safety risks. The authors should describe how they avoided releasing unsafe images.
        \item We recognize that providing effective safeguards is challenging, and many papers do not require this, but we encourage authors to take this into account and make a best faith effort.
    \end{itemize}

\item {\bf Licenses for existing assets}
    \item[] Question: Are the creators or original owners of assets (e.g., code, data, models), used in the paper, properly credited and are the license and terms of use explicitly mentioned and properly respected?
    \item[] Answer: \answerYes{} % Replace by \answerYes{}, \answerNo{}, or \answerNA{}.
    \item[] Justification: The models, datasets, evaluation tools, versions, and applicable licenses or terms of use are summarized in Table 6.
    \item[] Guidelines:
    \begin{itemize}
        \item The answer \answerNA{} means that the paper does not use existing assets.
        \item The authors should cite the original paper that produced the code package or dataset.
        \item The authors should state which version of the asset is used and, if possible, include a URL.
        \item The name of the license (e.g., CC-BY 4.0) should be included for each asset.
        \item For scraped data from a particular source (e.g., website), the copyright and terms of service of that source should be provided.
        \item If assets are released, the license, copyright information, and terms of use in the package should be provided. For popular datasets, \url{paperswithcode.com/datasets} has curated licenses for some datasets. Their licensing guide can help determine the license of a dataset.
        \item For existing datasets that are re-packaged, both the original license and the license of the derived asset (if it has changed) should be provided.
        \item If this information is not available online, the authors are encouraged to reach out to the asset's creators.
    \end{itemize}

\item {\bf New assets}
    \item[] Question: Are new assets introduced in the paper well documented and is the documentation provided alongside the assets?
    \item[] Answer: \answerNA{} % Replace by \answerYes{}, \answerNo{}, or \answerNA{}.
    \item[] Justification: No new assets are released with this submission.
    \item[] Guidelines:
    \begin{itemize}
        \item The answer \answerNA{} means that the paper does not release new assets.
        \item Researchers should communicate the details of the dataset\slash code\slash model as part of their submissions via structured templates. This includes details about training, license, limitations, etc. 
        \item The paper should discuss whether and how consent was obtained from people whose asset is used.
        \item At submission time, remember to anonymize your assets (if applicable). You can either create an anonymized URL or include an anonymized zip file.
    \end{itemize}

\item {\bf Crowdsourcing and research with human subjects}
    \item[] Question: For crowdsourcing experiments and research with human subjects, does the paper include the full text of instructions given to participants and screenshots, if applicable, as well as details about compensation (if any)? 
    \item[] Answer: \answerNA{} % Replace by \answerYes{}, \answerNo{}, or \answerNA{}.
    \item[] Justification: The paper does not involve crowdsourcing nor research with human subjects.
    \item[] Guidelines:
    \begin{itemize}
        \item The answer \answerNA{} means that the paper does not involve crowdsourcing nor research with human subjects.
        \item Including this information in the supplemental material is fine, but if the main contribution of the paper involves human subjects, then as much detail as possible should be included in the main paper. 
        \item According to the NeurIPS Code of Ethics, workers involved in data collection, curation, or other labor should be paid at least the minimum wage in the country of the data collector. 
    \end{itemize}

\item {\bf Institutional review board (IRB) approvals or equivalent for research with human subjects}
    \item[] Question: Does the paper describe potential risks incurred by study participants, whether such risks were disclosed to the subjects, and whether Institutional Review Board (IRB) approvals (or an equivalent approval/review based on the requirements of your country or institution) were obtained?
    \item[] Answer: \answerNA{} % Replace by \answerYes{}, \answerNo{}, or \answerNA{}.
    \item[] Justification: The paper does not involve crowdsourcing or research with human subjects.
    \item[] Guidelines:
    \begin{itemize}
        \item The answer \answerNA{} means that the paper does not involve crowdsourcing nor research with human subjects.
        \item Depending on the country in which research is conducted, IRB approval (or equivalent) may be required for any human subjects research. If you obtained IRB approval, you should clearly state this in the paper. 
        \item We recognize that the procedures for this may vary significantly between institutions and locations, and we expect authors to adhere to the NeurIPS Code of Ethics and the guidelines for their institution. 
        \item For initial submissions, do not include any information that would break anonymity (if applicable), such as the institution conducting the review.
    \end{itemize}

\item {\bf Declaration of LLM usage}
    \item[] Question: Does the paper describe the usage of LLMs if it is an important, original, or non-standard component of the core methods in this research? Note that if the LLM is used only for writing, editing, or formatting purposes and does \emph{not} impact the core methodology, scientific rigor, or originality of the research, declaration is not required.
    %this research? 
    \item[] Answer: \answerYes{} % Replace by \answerYes{}, \answerNo{}, or \answerNA{}.
    \item[] Justification: We use the externally provided gpt-4-1106-preview judge model solely for the standard AlpacaEval 2 utility evaluation. It is not part of the proposed validity-aware evaluation protocol or the merging methods, but it affects one reported downstream-utility metric.
    \item[] Guidelines:
    \begin{itemize}
        \item The answer \answerNA{} means that the core method development in this research does not involve LLMs as any important, original, or non-standard components.
        \item Please refer to our LLM policy in the NeurIPS handbook for what should or should not be described.
    \end{itemize}

\end{enumerate}





\end{document}